% Asterism agent_workflow — DeepMind co-math 教授簡介稿
% xelatex 編譯。中文 PMingLiU + Microsoft JhengHei。
\documentclass[11pt]{article}
\usepackage[a4paper, margin=2cm]{geometry}
\usepackage{xeCJK}
% Body: serif (LaTeX convention). Latin = CM Roman (default), CJK = PMingLiU.
\setCJKmainfont[AutoFakeBold=2.5, AutoFakeSlant=0.18]{PMingLiU}
% Figure sans (used in agent boxes for visual distinction): Microsoft JhengHei.
\setCJKsansfont{Microsoft JhengHei}
\setCJKmonofont{Microsoft JhengHei}
\usepackage{tikz}
\usetikzlibrary{positioning, arrows.meta, shapes.geometric, fit, backgrounds, calc}
\usepackage{booktabs, array}
\usepackage{enumitem}
\setlist{topsep=0.3em, itemsep=0.15em, parsep=0pt}

\setlength{\parindent}{0em}
\setlength{\parskip}{0.5em}

% ─── shared TikZ styles ───────────────────────────────────────────────
% Font scheme (uniform across all figures):
%   - swimlane / role headers: \sffamily\footnotesize
%   - node content (agent / state / hex / strat / file / db): \sffamily\scriptsize
%   - filenames in mono boxes: \ttfamily\scriptsize
%   - edge labels: \sffamily\scriptsize
\tikzset{
  agent/.style ={draw, rounded corners=2pt, inner sep=4pt,
                 font=\sffamily\scriptsize, fill=white, minimum height=0.6cm},
  human/.style ={draw, double, inner sep=3pt,
                 font=\ttfamily\scriptsize, fill=white},
  ffile/.style ={draw, inner sep=3pt,
                 font=\ttfamily\scriptsize, fill=white},
  db/.style    ={cylinder, draw, shape border rotate=90, aspect=0.35,
                 minimum width=1.1cm, minimum height=0.8cm,
                 font=\ttfamily\scriptsize, fill=white},
  hex/.style   ={draw, regular polygon, regular polygon sides=6,
                 inner sep=1pt, font=\sffamily\scriptsize, fill=white},
  st/.style    ={draw, inner sep=3pt, font=\sffamily\scriptsize,
                 fill=white, minimum height=0.6cm},
  state/.style ={draw, rounded corners=4pt, inner sep=3pt,
                 font=\sffamily\scriptsize, fill=white, minimum height=0.55cm},
  write/.style ={->, >=Stealth, semithick},
  read/.style  ={->, >=Stealth, semithick, dashed},
  arr/.style   ={->, >=Stealth, semithick},
  lbl/.style   ={font=\sffamily\scriptsize, inner sep=1.5pt, fill=white,
                 sloped=false},
}

\title{Asterism 系統介紹：自動化定理證明的多代理協作}
\author{}
\date{}

\begin{document}

\maketitle

本文件介紹 Asterism 框架的整體架構與運作流程。讀者預設熟悉數學與 Lean 4、但無工程背景。文中依主題搭配示意圖、表格或純文字說明。

「代理」（agent）指一次大型語言模型的呼叫；「框架」指環繞代理運作的調度、驗證、狀態管理基礎設施。

% =====================================================================
\section*{§1 \quad 動機}
% =====================================================================

直接請大型語言模型對一條 Lean 4 定理產出完整證明，往往失敗。常見失敗有三類：幻覺（援引不存在的引理）、無回饋（寫完後不知自己錯）、多步推理崩潰（單次呼叫的注意力無法處理數十步分解與串接）。

Asterism 的設計針對此三點：

\begin{itemize}
  \item \textbf{以 Lean 4 內核為事實依據}：任何證明步驟皆需通過型別檢查，杜絕幻覺。
  \item \textbf{透過 LSP 即時回饋}：代理寫一步即見編譯結果，無需等冷啟 \texttt{lake build}。
  \item \textbf{將大目標分派多個小代理}：每個代理只處理單一職責，避開單次呼叫的多步推理瓶頸。
\end{itemize}

% =====================================================================
\section*{§2 \quad 系統總覽：daemon 與 worker}
% =====================================================================

Asterism 在一次完整執行中由兩種角色構成：常駐 daemon 主導調度，按需 spawn 多個短命 worker 進行 LLM 呼叫。下圖顯示典型生命週期。

\begin{center}
\begin{tikzpicture}[
  header/.style={draw, fill=white, font=\sffamily\footnotesize,
                 inner sep=4pt, minimum height=0.7cm, align=center},
  note/.style ={draw, fill=white, font=\sffamily\scriptsize,
                inner sep=4pt, align=center},
  cond/.style ={draw, dashed, fill=white, font=\sffamily\scriptsize\itshape,
                inner sep=3pt},
  loopf/.style={draw, fill=white, font=\sffamily\scriptsize\itshape,
                inner sep=2pt},
  msg/.style  ={font=\sffamily\scriptsize, midway, above, fill=white,
                inner sep=2pt},
  ret/.style  ={font=\sffamily\scriptsize, midway, above, fill=white,
                inner sep=2pt},
  lane/.style ={thin, gray!60},
]
  % --- swimlane headers ---
  \node[header] (uH) at (0,    5)   {使用者};
  \node[header] (dH) at (3.5,  5)   {daemon\\(常駐)};
  \node[header] (wH) at (8,    5)   {worker\\(短命)};
  \node[header] (lH) at (12,   5)   {Lean 內核};
  \node[header] (uF) at (0,   -5.2) {使用者};
  \node[header] (dF) at (3.5, -5.2) {daemon};
  \node[header] (wF) at (8,   -5.2) {worker};
  \node[header] (lF) at (12,  -5.2) {Lean 內核};
  \draw[lane] (uH.south) -- (uF.north);
  \draw[lane] (dH.south) -- (dF.north);
  \draw[lane] (wH.south) -- (wF.north);
  \draw[lane] (lH.south) -- (lF.north);
  % --- entry: cli run ---
  \draw[arr] (0, 4) -- (3.5, 4) node[msg]{cli run};

  % --- loop fragment (UML-style): rectangle around the iterated section ---
  \draw[thin] (1.6, 3.5) rectangle (13.2, -3.0);
  \node[loopf, anchor=north west] at (1.6, 3.5) {loop \upshape\sffamily{[未達根目標 \texttt{proved}]}};

  % --- messages inside loop ---
  \node[note] at (3.5, 2.7) {主迴圈：每輪掃描 +\\worker 完成觸發};
  \draw[arr] (3.5, 1.8) -- (8, 1.8) node[msg]{spawn worker (一次 LLM 呼叫)};
  \draw[arr] (8, 1.0) -- (12, 1.0) node[msg]{即時查詢編譯結果};
  \draw[arr, dashed] (12, 0.3) -- (8, 0.3) node[ret]{錯誤 / 目前 goal};
  \draw[arr, dashed] (8, -0.5) -- (3.5, -0.5) node[ret]{提交結果};
  \node[note] at (3.5, -1.4) {更新資料庫 + 串聯處理};
  \node[cond, anchor=west] at (4.5, -2.5) {跳出條件：所有根目標 \texttt{proved}};

  % --- exit ---
  \draw[arr, dashed] (3.5, -4.0) -- (0, -4.0) node[ret]{daemon 退出};
\end{tikzpicture}
\end{center}

\begin{itemize}
  \item \textbf{daemon}：常駐 Python 進程、負責調度迴圈。每次掃描判定「下一步該派哪個 worker」、收到 worker 結果後更新資料庫並決定串聯效應（例如某個子目標證出後是否解鎖父目標）。
  \item \textbf{worker}：每個 worker 即一次 LLM 呼叫的容器、做完即退出。沙盒目錄為 \texttt{.attempts/<uuid>/}，預設工作池容量為 4（即同一時間最多 4 個 worker 並行）。
  \item \textbf{Lean 內核}：daemon 在背景駐留一個 Lean 伺服器、提供 worker 寫每一步證明時的即時編譯回饋（無需冷啟 \texttt{lake build} 的 30–60 秒成本）。
\end{itemize}

整個生命週期由 daemon 主導；worker 僅回傳結果即退出，直至所有納入該次執行的問題之根目標皆達 \texttt{proved}，daemon 才終止。

% =====================================================================
\section*{§3 \quad 檔案的所有權與讀寫權限}
% =====================================================================

各 pipeline 對檔案的動作分為兩個階段。圖中實線箭頭表示寫動作、虛線箭頭表示讀動作；箭頭旁的小字標註觸發時機。雙框節點為人手檔（pipeline 只讀）、單框節點為框架檔。

\subsection*{A.\quad 啟動階段（一次性）}

\begin{center}
\begin{tikzpicture}
  \node[agent] (user)  at (0,  1.5) {使用者};
  \node[human] (manif) at (4,  2.4) {Manifest.md};
  \node[human] (defs)  at (4,  0.6) {Defs.lean};
  \node[agent] (cli)   at (8,  1.5) {cli init};
  \node[ffile] (root)  at (12, 2.4) {Root.lean};
  \node[db]    (db)    at (12, 0.6) {asterism.db};

  \draw[write] (user) to[bend left=10]  node[lbl, above]{建立 / 編輯} (manif);
  \draw[write] (user) to[bend right=10] node[lbl, below]{建立 / 編輯} (defs);
  \draw[read]  (manif) to[bend left=10] node[lbl, above]{讀} (cli);
  \draw[read]  (defs)  to[bend right=10] node[lbl, below]{讀} (cli);
  \draw[write] (cli) to[bend left=10]   node[lbl, above]{寫含 \texttt{sorry} 的占位檔} (root);
  \draw[write] (cli) to[bend right=10]  node[lbl, below]{寫 schema} (db);
\end{tikzpicture}
\end{center}

使用者寫好問題定義（\texttt{Manifest.md} 框出 root theorem 與 axioms whitelist；\texttt{Defs.lean} 補必要的 def / structure）後 \texttt{cli init} 一次性把 \texttt{Root.lean}（含 \texttt{sorry} 占位）與 \texttt{asterism.db}（schema 建好、root goal 入庫）寫出。其後該問題的人手檔不再被任何 pipeline 寫。

\subsection*{B1.\quad 運作階段 — 正常讀寫流（每 spawn）}

\begin{center}
\begin{tikzpicture}[
  ephem/.style ={draw, dashed, inner sep=3pt,
                 font=\ttfamily\scriptsize, fill=white},
]
  \node[human] (manif)   at (0,   5.5) {Manifest.md};
  \node[ffile] (lessons) at (0,   3.7) {LESSONS.md};
  \node[ephem] (ctx)     at (3.5, 4.6) {Context.md};
  \node[agent] (strat)   at (7,   6.8) {Strategist};
  \node[agent] (back)    at (7,   5.0) {Backward};
  \node[agent] (build)   at (7,   3.2) {Builder};
  \node[agent] (fwd)     at (7,   0.6) {Forward};
  \node[ephem] (att)     at (10.5, 1.9) {.attempts/<pid>/};
  \node[db]    (db)      at (13.5, 6.8) {asterism.db};
  \node[ffile] (strat_f) at (13.5, 5.0) {\_strategy\_s*.lean};
  \node[ffile] (proofs)  at (13.5, 3.2) {proofs/L\_*.lean};

  \draw[write] (manif)   to[bend left=4]  node[lbl, above, pos=0.6]{spawn 前 框架編譯} (ctx);
  \draw[write] (lessons) to[bend right=4] node[lbl, below, pos=0.6]{(Backward/Builder 才含)} (ctx);
  \draw[read]  (ctx) to[bend left=4]  node[lbl, above, pos=0.6]{spawn 時 讀} (strat);
  \draw[read]  (ctx) -- (back);
  \draw[read]  (ctx) to[bend right=2] (build);
  \draw[read]  (ctx) to[bend right=10] (fwd);
  \draw[write] (back)  to[bend left=8]  node[lbl, above, pos=0.45]{分解輸出} (att);
  \draw[write] (build) to[bend left=4]  node[lbl, above]{證明輸出} (att);
  \draw[write] (fwd)   to[bend right=15] node[lbl, below, pos=0.55]{新引理輸出} (att);
  \draw[write] (att) to[bend left=14] node[lbl, above, pos=0.5]{框架 promote} (strat_f);
  \draw[write] (att) to[bend right=8] node[lbl, below, pos=0.5]{框架 promote} (proofs);
  \draw[write] (strat) -- node[lbl, below]{commit decision + directive} (db);
\end{tikzpicture}
\end{center}

\texttt{Context.md} 與 \texttt{.attempts/<pid>/} 皆為 per-spawn ephemeral（虛框）：框架在 spawn 開始前寫入、worker 退出時 \texttt{WorkArea.\_\_exit\_\_} 整目錄 \texttt{rmtree}。Agent 從不直接讀 \texttt{Manifest.md} 與 \texttt{LESSONS.md}，一律讀框架已編譯好的當輪 \texttt{Context.md}；agent 輸出也是先寫 sandbox、框架驗證通過後才 promote 到永久檔，驗證失敗則 sandbox 整目錄 rmtree、永久檔不被汙染。

Context 內容因 agent 種類而異：Strategist 與 Forward 額外注入 \texttt{TREE.md}（problem snapshot），Strategist 另含 db 狀態切片（active goals、近期決策）；\texttt{Defs.lean} 不進 Context，agent 透過 Lean LSP 在 elaboration 時間取用。

\subsection*{B2.\quad 運作階段 — 框架托管（每 tick／pipeline 收尾／超時）}

\begin{center}
\begin{tikzpicture}[
  ephem/.style ={draw, dashed, inner sep=3pt,
                 font=\ttfamily\scriptsize, fill=white},
]
  \node[ephem] (att)    at (0,   3.5) {.attempts/<pid>/};
  \node[ffile] (drafts) at (4.5, 3.5) {.drafts/};
  \node[ephem] (ctxN)   at (9.5, 3.5) {Context.md \textit{(下次 spawn)}};
  \draw[write] (att)    -- node[lbl, above]{超時 postmortem 寫筆記} (drafts);
  \draw[write] (drafts) -- node[lbl, above]{重派時 注入} (ctxN);

  \node[agent] (refl)    at (0,   1.8) {reflection};
  \node[ffile] (lessons) at (4.5, 1.8) {LESSONS.md};
  \draw[write] (refl) -- node[lbl, above]{pipeline 收尾 寫追加} (lessons);

  \node[agent] (verify)  at (0,  -0.8) {verify};
  \node[db]    (db)      at (5,   0.5) {asterism.db};
  \node[ffile] (root)    at (5,  -2.4) {Root.lean};
  \node[ffile] (proofs)  at (9.5,-0.8) {proofs/L\_*.lean};
  \draw[write] (verify) to[bend left=10] node[lbl, above, pos=0.55]{每輪 寫狀態} (db);
  \draw[write] (verify) to[bend right=10] node[lbl, above, pos=0.65]{每輪 reconcile (Root)} (root);
  \draw[write] (verify) -- node[lbl, above, pos=0.6]{每輪 reconcile (L\_*)} (proofs);
\end{tikzpicture}
\end{center}

B2 的三條流程都不涉及 LLM：postmortem 寫筆記由 daemon 處理（agent 不參與）；reflection 由框架在 pipeline 收尾時自動 spawn 短呼叫；verify housekeeping 在 dispatcher 主迴圈每輪執行（純 Python、不佔 worker pool）。

人手只動 \texttt{Manifest.md} 與 \texttt{Defs.lean}，其餘檔案皆框架擁有；手改會使 \texttt{asterism.db}（SoT）與檔案系統失同步。代理欲修正人手檔須走 \texttt{RequestUserAmend} 暫停問題等待人工確認。

% =====================================================================
\section*{§4 \quad 四個代理}
% =====================================================================

Asterism 將證明工作分派給四個專責代理。各代理對應一次大型語言模型呼叫，輸入輸出與時間預算各異：

\begin{center}
\renewcommand{\arraystretch}{1.3}
\begin{tabular}{@{}p{2.2cm} p{6cm} p{6cm}@{}}
\toprule
\textbf{代理} & \textbf{輸入} & \textbf{輸出} \\
\midrule
Strategist & 全域進度視圖、最近決策歷史、TREE 快照、注入批次結果、\texttt{pending\_strategist\_review} 待裁決清單、Manifest meta & 一筆或多筆決策（注入、放行、確認擱置等可組成 batch） \\
\midrule
Backward & 一個待拆目標 + 父策略提示 & 0 至 N 個子目標 + 結構性組合命令（0 個 = 葉節點直接證） \\
\midrule
Builder & 一個原子目標 & 完整證明（可通過型別檢查與 axiom 探針） \\
\midrule
Forward & Strategist 給的 brief & 一條新的 theorem / def / structure / class 檔案 \\
\bottomrule
\end{tabular}
\end{center}

四者形成分層分工：Strategist 規劃調度、Backward 結構分解、Builder 完成原子證明、Forward 補前置引理。

% =====================================================================
\section*{§5 \quad 端到端流程}
% =====================================================================

以下時序圖為一條典型問題從提交到證完的完整過程。

\begin{center}
\begin{tikzpicture}[
  header/.style={draw, fill=white, font=\sffamily\footnotesize,
                 inner sep=4pt, minimum height=0.7cm, align=center,
                 minimum width=1.7cm},
  note/.style ={draw, fill=white, font=\sffamily\scriptsize,
                inner sep=3pt, align=center},
  phase/.style={draw, dashed, fill=white, font=\sffamily\scriptsize\itshape,
                inner xsep=6pt, inner ysep=2pt},
  loopf/.style={draw, fill=white, font=\sffamily\scriptsize\itshape,
                inner sep=2pt},
  msg/.style  ={font=\sffamily\scriptsize, midway, above, fill=white,
                inner sep=1.5pt},
  lane/.style ={thin, gray!60},
]
  % --- swimlane headers ---
  \node[header] (uH)  at (0,     11)   {使用者};
  \node[header] (sH)  at (2.6,   11)   {Strategist};
  \node[header] (dH)  at (5.2,   11)   {Dispatcher};
  \node[header] (fH)  at (7.8,   11)   {Forward};
  \node[header] (bH)  at (10.4,  11)   {Backward};
  \node[header] (bdH) at (13,    11)   {Builder};
  % --- footers ---
  \node[header] (uF)  at (0,    -8.5)  {使用者};
  \node[header] (sF)  at (2.6,  -8.5)  {Strategist};
  \node[header] (dF)  at (5.2,  -8.5)  {Dispatcher};
  \node[header] (fF)  at (7.8,  -8.5)  {Forward};
  \node[header] (bF)  at (10.4, -8.5)  {Backward};
  \node[header] (bdF) at (13,   -8.5)  {Builder};
  % --- lifelines ---
  \foreach \x in {0, 2.6, 5.2, 7.8, 10.4, 13} \draw[lane] (\x, 10.6) -- (\x, -8.1);

  % --- Phase 1: bootstrap + first_launch ---
  \node[phase, anchor=west] at (-1, 10.2) {Phase 1 — bootstrap};
  \draw[arr] (0,    9.6) -- (5.2,  9.6) node[msg]{寫 Manifest + Defs};
  \draw[arr] (5.2,  8.8) -- (2.6,  8.8) node[msg]{\texttt{first\_launch} 觸發};
  \node[note] at (2.6, 8.0) {Grep mathlib +\\EmitDirective};
  \draw[arr, dashed] (2.6, 7.0) -- (5.2, 7.0) node[msg]{\texttt{Inject(brief=X)} 決策};
  \draw[arr] (5.2,  6.2) -- (7.8,  6.2) node[msg]{下一輪 BFS enqueue};
  \draw[arr, dashed] (7.8, 5.4) -- (5.2, 5.4) node[msg]{X 已證 (cascade)};
  \draw[arr] (5.2,  4.6) -- (2.6,  4.6) node[msg]{\texttt{inject\_batch\_done} 觸發};
  \draw[arr, dashed] (2.6, 3.8) -- (5.2, 3.8) node[msg]{Reopen 根目標};

  % --- Phase 2: BFS dispatch + Backward + parallel Builder ---
  % loop fragment: dispatcher iterates BFS until queue empty / all proved
  \node[phase, anchor=west] at (-1, 3.0) {Phase 2 — BFS（dispatcher 反覆派工）};
  \draw[thin] (1.2, 2.7) rectangle (13.6, -0.4);
  \node[loopf, anchor=north west] at (1.2, 2.7)
        {loop \upshape\sffamily{[queue 不空]}};
  \draw[arr] (5.2,  2.0) -- (10.4, 2.0) node[msg]{BFS enqueue};
  \draw[arr, dashed] (10.4, 1.3) -- (5.2, 1.3) node[msg]{拆出 g1, g2, g3};
  \draw[arr] (5.2,  0.6) -- (13,   0.6) node[msg]{三個並行 enqueue};
  \draw[arr, dashed] (13, -0.1) -- (5.2, -0.1) node[msg]{g1, g3 proved；g2 \texttt{return\_to\_parent}};

  % --- Phase 3: pending_review + reopen + success ---
  \node[phase, anchor=west] at (-1, -0.8) {Phase 3 — 失敗恢復};
  \draw[arr] (5.2, -1.4) -- (2.6, -1.4) node[msg]{\texttt{pending\_review} 觸發 (g2)};
  \draw[arr, dashed] (2.6, -2.2) -- (5.2, -2.2) node[msg]{Reopen g2 換角度};
  \draw[arr] (5.2, -3.0) -- (10.4, -3.0) node[msg]{BFS enqueue g2};
  \draw[arr, dashed] (10.4, -3.8) -- (5.2, -3.8) node[msg]{新拆 g4–g7};
  \draw[arr] (5.2, -4.6) -- (13,  -4.6) node[msg]{並行 enqueue};
  \draw[arr, dashed] (13, -5.4) -- (5.2, -5.4) node[msg]{全部 proved};
  \node[note] at (5.2, -6.2) {root 連鎖 \texttt{proved}};
  \draw[arr, dashed] (5.2, -7.0) -- (0, -7.0) node[msg]{daemon 退出};
\end{tikzpicture}
\end{center}

每個步驟皆由框架自動觸發、無需人工介入。Strategist 只寫決策（Inject / Reopen / ConfirmShelve 等），不直接派工；Dispatcher 是唯一的「行動者」，在下一輪 BFS 中根據 Strategist 寫入的決策與當前 goal 狀態 enqueue worker。失敗（如 g2 的 \texttt{return\_to\_parent} 宣告）為一等公民、由 Strategist 介入後重新規劃。

% =====================================================================
\section*{§6 \quad AND/OR 圖：證明搜索的資料結構}
% =====================================================================

證明搜索的中心資料結構為 AND/OR 圖。OR 節點為目標、AND 節點為策略：

\begin{center}
\begin{tikzpicture}[
  goal/.style={hex, minimum size=1cm},
  proved/.style={hex, minimum size=1cm, fill=black!10},
  dead/.style={hex, minimum size=1cm, dashed, very thick},
  strat/.style={st, minimum width=2cm, minimum height=0.7cm, align=center},
  deadstrat/.style={st, minimum width=2cm, minimum height=0.7cm, dashed, very thick, align=center},
  edge/.style={-, semithick},
  deadedge/.style={-, semithick, dashed},
]
  \node[goal] (root) at (5, 5) {main};
  \node[strat] (s1) at (2.5, 3) {策略 s1\\歸納分解};
  \node[deadstrat] (s2) at (7.5, 3) {策略 s2\\反證法 (死)};
  \node[proved] (g1) at (0.5, 0.5) {g1 \checkmark};
  \node[goal] (g2) at (2.5, 0.5) {g2};
  \node[goal] (g3) at (4.5, 0.5) {g3};
  \node[dead] (g4) at (6.5, 0.5) {g4};
  \node[dead] (g5) at (8.5, 0.5) {g5};

  \draw[edge] (root) -- node[lbl, left, pos=0.5]{擇一即可} (s1);
  \draw[deadedge] (root) -- (s2);
  \draw[edge] (s1) -- node[lbl, left, pos=0.6]{皆需證} (g1);
  \draw[edge] (s1) -- (g2);
  \draw[edge] (s1) -- (g3);
  \draw[deadedge] (s2) -- (g4);
  \draw[deadedge] (s2) -- (g5);
\end{tikzpicture}
\end{center}

\begin{itemize}
  \item 六邊形節點為\textbf{目標}（OR）；任一掛在其下的策略證出即視為證明。
  \item 矩形節點為\textbf{策略}（AND）；所有掛在其下的子目標皆需證出，策略才算成功。
  \item 實線 + 細邊：活路徑。粗虛線：已死路徑（圖中策略 s2 與其子目標 g4、g5）。淺灰底：已證子目標（圖中 g1）。
  \item 同一目標下多個策略\textbf{序列}嘗試而非並行（passive OR、cap = 1）：當前活策略死透後，框架才允許 Backward 為該目標生第二條策略。早期 eager fanout 在強模型下實證為純粹浪費 token。
  \item 同一策略下的兄弟子目標可並行（受 \texttt{dispatch.pool} 上限約束、預設 4）。
  \item 框架以廣度優先搜尋掃描所有「存活鏈通達根」的 \texttt{open} 目標、按優先級派工；死路徑不影響其他路徑繼續推進。
\end{itemize}

% =====================================================================
\section*{§7 \quad 目標狀態機}
% =====================================================================

每個目標在資料庫中有一個狀態欄位，所有狀態轉移由框架統一管理：

\begin{center}
\begin{tikzpicture}[
  s/.style    ={state, minimum width=2cm, minimum height=0.75cm},
  init/.style ={circle, fill=black, inner sep=2pt},
  termmark/.style={state, double, minimum width=2cm, minimum height=0.75cm},
]
  % --- initial dot + main horizontal axis at y=0 ---
  \node[init] (start) at (-1.8, 0)  {};
  \node[s] (frozen) at (0,    0)    {frozen};
  \node[s] (open)   at (4,    0)    {open};
  \node[s] (atte)   at (9,    0)    {attempting};
  \node[s] (psr)    at (14,   0)    {pending\_review};
  \node[s, dashed] (shelv) at (14, -4) {shelved};
  % --- proved above attempting (success goes up, rollback comes down) ---
  \node[termmark] (proved) at (9, 4) {proved};
  % --- hard terminals: dead well left of attempting, disproved well right ---
  \node[termmark] (dead)  at (4.5,  -4) {dead};
  \node[termmark] (dispr) at (10,   -4) {disproved};

  % ─── forward axis (y=0) ─────────────────────────────────────────
  \draw[arr] (start)  -- node[lbl, above]{cli init} (frozen);
  \draw[arr] (frozen) -- node[lbl, above]{first\_launch Reopen} (open);
  \draw[arr] (open)   -- node[lbl, above]{BFS 派工} (atte);
  \draw[arr] (atte) -- node[lbl, above]{agent shelve / 達上限} (psr);

  % ─── attempting ↔ proved (parallel arrows, offset to avoid overlap) ──
  \draw[arr] ([xshift=-0.3cm]atte.north) -- node[lbl, left]{所有 sub-goal 證出} ([xshift=-0.3cm]proved.south);
  \draw[arr, dashed] ([xshift=0.3cm]proved.south) -- node[lbl, right]{rollback（sorryAx 偵測，罕見）} ([xshift=0.3cm]atte.north);

  % ─── attempting → hard terminals (below, label on outer side) ────
  \draw[arr] (atte) -- node[lbl, left, pos=0.6]{父策略 cascade} (dead);
  \draw[arr] (atte) -- node[lbl, right, pos=0.6]{agent 給反例} (dispr);

  % ─── retry loop: attempting → open (clean low-bend below axis) ──
  \draw[arr] (atte) to[bend left=18]
      node[lbl, below, pos=0.5]{retry（\texttt{attempts++} < threshold）} (open);

  % ─── pending_review → open (上方大弧 Reopen) ──────────────────
  \draw[arr] (psr) to[bend right=22]
      node[lbl, above, pos=0.5]{Strategist Reopen} (open);

  % ─── pending_review → shelved (vertical, ConfirmShelve) ──────
  \draw[arr] (psr) -- node[lbl, right]{ConfirmShelve} (shelv);

  % ─── shelved → open (軟可逆 Reopen，繞最下方避開 dead/disproved) ─
  \draw[arr, rounded corners=8pt] (shelv) |- ($(open.south)+(0,-6)$) -| (open.south);
  \node[lbl, anchor=south] at (9, -6.2) {Strategist Reopen（軟可逆）};

\end{tikzpicture}
\end{center}

關鍵設計：

\begin{itemize}
  \item \texttt{shelved} 與 \texttt{dead} 區分軟硬終結。\texttt{shelved} 仍可被 Strategist 透過 Reopen 復活（例如後續注入的引理改變了可證性）。\texttt{dead} 為母策略整個被否決所致，本目標於該策略內不再嘗試。
  \item \texttt{pending\_strategist\_review} 為過渡狀態，代理主動聲明無解後置入，等候 Strategist 裁決；本身不終結，故串聯不向上傳播。
  \item \texttt{proved} 非絕對終態：root integrity gate 在每次 root 完工時對 \texttt{Root.lean} 做 axiom probe，若偵測到 \texttt{sorryAx} 殘留即走 \texttt{bisect\_sorryax\_source} + \texttt{rollback\_cascade\_chain}，元凶策略撤回、下游 goal 回 \texttt{attempting}。實證自 verify-collapse 以來 41+ 輪 0 次觸發，gate 設計用於極端 corner case 而非常規路徑。
\end{itemize}

% =====================================================================
\section*{§8 \quad 失敗處理流程}
% =====================================================================

\begin{center}
\begin{tikzpicture}[
  proc/.style ={draw, rounded corners=2pt, inner sep=4pt,
                font=\sffamily\scriptsize, fill=white,
                minimum height=0.6cm, align=center},
  decis/.style={draw, diamond, aspect=2.2, inner sep=2pt,
                font=\sffamily\scriptsize, fill=white},
  term/.style ={draw, inner sep=4pt, font=\sffamily\scriptsize,
                fill=white, align=center},
  cyl/.style  ={cylinder, draw, shape border rotate=90, aspect=0.4,
                minimum width=1.6cm, minimum height=0.7cm,
                font=\ttfamily\scriptsize, fill=white},
  elabel/.style={font=\sffamily\scriptsize, fill=white, inner sep=2pt},
]
  % ─── Top: worker + outcome decision ─────────────────────────────
  \node[proc]  (work) at (8, 9.5) {Worker 執行};
  \node[decis] (res)  at (8, 8.2) {結果?};
  \draw[arr] (work) -- (res);

  % ─── 4 first-level outcome columns ──────────────────────────────
  \node[proc] (ok)    at (1.5,  6.8) {推進至 workspace\\+ axiom 探針};
  \node[decis] (kind) at (5.5,  6.8) {哪種?};
  \node[proc] (kill)  at (10.5, 6.8) {Watchdog SIGKILL};
  \node[proc] (infra) at (14.5, 6.8) {infra 退避 30\,s\\冷卻};
  \draw[arr] (res) -- node[elabel, above, pos=0.7]{成功} (ok);
  \draw[arr] (res) -- node[elabel, right, pos=0.55]{寫 decline} (kind);
  \draw[arr] (res) -- node[elabel, above, pos=0.7]{超時} (kill);
  \draw[arr] (res) -- node[elabel, above, pos=0.8]{spawn 失敗} (infra);

  % ─── Branch 1: 成功 → reflection → LESSONS ─────────────────────
  \node[proc] (refl) at (1.5, 4.7) {reflection 寫追加};
  \node[cyl]  (less) at (1.5, 2.8) {LESSONS.md};
  \draw[arr] (ok)   -- (refl);
  \draw[arr] (refl) -- (less);

  % ─── Branch 2: decline → 3 sub-cases (vertical stack below kind) ─
  \node[term] (dispr) at (5.5, 5.1) {\texttt{unprovable} \;$\to$\; disproved};
  \node[term] (dead)  at (5.5, 3.9) {\texttt{return\_to\_parent} \;$\to$\; dead};
  \node[term] (psr)   at (5.5, 2.7) {\texttt{shelve} \;$\to$\; pending\_review};
  \draw[arr] (kind) -- (dispr);
  \draw[arr] (dispr) -- (dead);
  \draw[arr] (dead)  -- (psr);

  \node[term] (casc) at (5.5, 1.4)
        {dead 後續：Cascade（標記父策略 + \texttt{attempts++}）};
  \draw[arr, dashed] (dead.east) to[bend right=20] (casc.east);

  % ─── Branch 3: 超時 → postmortem → .drafts/ ─────────────────────
  \node[proc] (pm)     at (10.5, 4.9) {postmortem\\(180\,s, \texttt{--resume sid})};
  \node[cyl]  (drafts) at (10.5, 2.8) {.drafts/};
  \draw[arr] (kill) -- (pm);
  \draw[arr] (pm)   -- (drafts);

\end{tikzpicture}
\end{center}

四類失敗對應不同處理路徑。所有處理由框架統一執行、代理無自主決定權。

\begin{itemize}
  \item \textbf{decline 指令}：代理在沙盒中探索後判定無解、主動聲明，附理由與修正建議。
  \item \textbf{超時}：watchdog 偵測長時間無工具呼叫即 SIGKILL；接著 pipeline 自動以 \texttt{claude --resume <sid>} 從磁碟救回該 session、再 spawn 一輪短呼叫（180s 上限）請代理寫進度筆記。此筆記隨後注入新代理的 Context、使重派時不致從零開始。
  \item \textbf{spawn 失敗}：外部基礎設施問題（如配額耗盡），不計入嘗試次數。
  \item \textbf{成功}：reflection 短呼叫請代理回看本次經驗、寫入跨 spawn 累積的 \texttt{LESSONS.md}。
\end{itemize}

% =====================================================================
\section*{§9 \quad Strategist 的四種觸發條件}
% =====================================================================

Strategist 不主動工作，必待事件觸發。當前定義四種觸發類型，每種對應一份獨立的指引提示：

\begin{center}
\renewcommand{\arraystretch}{1.3}
\begin{tabular}{@{}p{3.3cm} p{6.5cm} p{4.5cm}@{}}
\toprule
\textbf{觸發} & \textbf{何時觸發} & \textbf{主要任務} \\
\midrule
\texttt{first\_launch} & 問題首次啟動（root 仍 \texttt{frozen} 且尚無任何 Strategist 決策） & 調研 mathlib、發布指導、派出前置 \\
\midrule
\texttt{routine} & 每小時定期 & 稽核分支、識別停滯、發現重造輪子 \\
\midrule
\texttt{pending\_review} & 子目標出現「待審」 & 裁決失敗：重派、改方向、確認放棄 \\
\midrule
\texttt{inject\_batch\_done} & 注入批次全部終結 & 評估結果、決定後續（多半 Reopen 被擱置目標） \\
\bottomrule
\end{tabular}
\end{center}

每個觸發對應一份獨立 prompt：代理被呼叫時只見該觸發的指引、不受其他模式干擾、注意力集中。觸發優先級由上而下：\texttt{inject\_batch\_done} 最高、\texttt{routine} 最低。

\texttt{routine} 為唯一無事件驅動的觸發、其職責為防止系統長時間停滯而無人察覺；當前 prompt 明確要求代理執行品質稽核清單、而非預設回應「無事可做」。

% =====================================================================
\section*{§10 \quad 工具與資料管道}
% =====================================================================

每個代理被呼叫時，框架提供一組工具與分層資料來源。

\subsection*{工具}

\begin{itemize}
  \item 檔案操作：\texttt{Read} / \texttt{Write} / \texttt{Edit} / \texttt{Grep}，標準命令列語意。
  \item 與 Lean 伺服器交互：\texttt{apply\_edit} / \texttt{goal\_at} / \texttt{errors\_at} / \texttt{validate\_file}，皆於 1 秒內回應，避免冷啟 \texttt{lake build} 的 30–60 秒成本。
  \item Mathlib 搜尋：\texttt{Grep} 適用於關鍵字、\texttt{loogle} 適用於以陳述形狀搜尋。
\end{itemize}

\subsection*{資料來源（三層）}

第一層為\textbf{任務簡報} \texttt{Context.md}：每次呼叫前框架即時編譯一份、即用即棄。內容因代理種類而異——Strategist 看全域進度視圖與最近決策歷史；Backward 與 Builder 看當前目標陳述、父策略提示、可引用的兄弟證明；Forward 看 Strategist 指定的注入 brief。

第二層為\textbf{跨 spawn 累積的經驗} \texttt{LESSONS.md}：由 reflection 階段在 worker 收場時寫入，典型內容如「mathlib 已有 X 引理、不必重造」「Y 路線是死路」「Z 型別陷阱」等戰術提示。每個代理冷啟時皆會於 Context 中看到當前內容。

第三層為\textbf{問題級常駐指引} Strategist directive：Strategist 透過 \texttt{EmitDirective} 寫入，作為該問題下所有後續代理的長期方向。每次更新覆寫前次內容；每個 worker 冷啟皆會看到最新版本。

資料的流動方向呈一鏈狀：使用者寫 Manifest → Strategist 讀取後寫 directive 與 LESSONS 一同進到 worker → worker 工作完畢由 reflection 寫回 LESSONS → 下一次代理冷啟時讀到。此鏈使知識跨多次嘗試累積、不致每次冷啟皆從零開始。

% =====================================================================
\section*{§11 \quad 現況與成果}
% =====================================================================

\textbf{已自動證出問題涵蓋：}

\begin{center}
\renewcommand{\arraystretch}{1.3}
\begin{tabular}{@{}p{2.5cm} p{12cm}@{}}
\toprule
\textbf{領域} & \textbf{代表問題} \\
\midrule
數論 & Wilson 定理、Pythagorean 三元組、若干 Minif2f 競賽題 \\
\midrule
拓樸 & $\pi_1(S^1)$ 基本群、Sylvester–Gallai 共線定理 \\
\midrule
線性代數 & Schur 三角化、SVD、極分解、QR 分解 \\
\midrule
分析 & 留數定理 \\
\bottomrule
\end{tabular}
\end{center}

\textbf{典型成本}：每個問題從提交到證完，平均耗時 2–8 小時、佔 LLM 呼叫 50–300 次。問題的內在複雜度（分解深度、引用 mathlib 引理數量、是否需要新工具）決定實際成本。

\textbf{框架自身}：累計 1141 項自動測試、涵蓋狀態機轉換、串聯規則、注入決策、復原邏輯等核心不變式。每次程式碼變更皆需通過全部測試。

\textbf{近期觀察到的失敗模式：}

\begin{itemize}
  \item \textbf{重造輪子}：深度極大的子目標常為 Backward 拆題時未識別到 mathlib 已有對應引理。已新增 Citation pass 流程。
  \item \textbf{觸發保護過嚴}：早期 Strategist 觸發互不覆蓋導致停滯。已加結構性停滯偵測。
  \item \textbf{代理注意力被稀釋}：LESSONS 累積過多時代理略讀。下一步規劃關鍵字標籤化。
\end{itemize}

% =====================================================================
\section*{§12 \quad 設計取捨}
% =====================================================================

\begin{center}
\renewcommand{\arraystretch}{1.3}
\begin{tabular}{@{}p{2.3cm} p{6.5cm} p{6cm}@{}}
\toprule
\textbf{取捨} & \textbf{我們的選擇} & \textbf{代價} \\
\midrule
驗證方式 & Lean 4 內核（唯一事實源） & 速度較慢、受限 mathlib 覆蓋範圍 \\
\midrule
狀態源 & SQLite 資料庫（檔案系統為衍生物） & 手動編輯狀態檔風險極高 \\
\midrule
代理結構 & 多代理分工（4 種角色） & 上下文連續性犧牲、換取注意力聚焦 \\
\midrule
失敗處理 & 代理可自主 decline + Strategist 接手 & 仰賴 Strategist 判斷品質 \\
\midrule
反思機制 & reflection 後寫 LESSONS、跨 spawn 累積 & LESSONS 過長時略讀風險 \\
\bottomrule
\end{tabular}
\end{center}

\textbf{取捨核心}：選擇 Lean 4 作為唯一事實依據，意味系統無需另設「評審代理」（如其他類似系統採用）——Lean 內核即評審。代價為代理寫的每一步皆需通過嚴格型別檢查、產出較慢。

\textbf{未解問題：}

\begin{itemize}
  \item \textbf{跨問題重用}：「Library」子模組尚未啟用，不同問題間無法共享已證引理。
  \item \textbf{長尾子目標}：少數深度極大的子目標耗時遠超預期。
  \item \textbf{經驗傳遞對齊}：LESSONS 與 directive 的代理利用率仍待提升。
\end{itemize}

下一階段規劃將成熟的引理自動晉升為 Library 條目，並引入關鍵字標籤化機制集中代理注意力於最相關的歷史經驗。

\end{document}
